
\question
Consider the food chain grass $\to$ grasshopper $\to$ mouse $\to$ snake $\to$ hawk. How much of the chemical energy produced by photosynthesis of the grass (100\%) is available to the hawk?

\begin{choices}
\correctchoice 0.01\%
\choice 10\%
\choice 1\%
\choice 0.1\%
\choice 60\%
\end{choices}

\question
Which of the following areas of study focuses on the exchange of energy and nutrients among individuals of different species?

\begin{choices}
\choice Behavioral Ecology 
\choice Evolutionary Ecology 
\correctchoice Ecosystem ecology. 
\choice Community Ecology 
\choice Population Ecology 
\end{choices}

\newpage

\question
A genetic change that caused a certain \textit{Hox} gene to be expressed along the tip of a vertebrate limb bud instead of farther back helped make possible the  evolution of the tetrapod limb. This type of change is an example of 
\begin{choices}
\choice the influence of environment on development.
\choice heterochrony.
\choice paedomorphosis.
\correctchoice a change in a toolkit gene that altered the organization of body parts.
\choice allometry.
\end{choices}

\question
A small population of white-footed mice has the same intrinsic rate of increase (\textit{r}) as a large population.  If everything else is equal, 

\begin{choices}
\choice the growth trajectories of the two populations will proceed in opposite directions. 
\choice the two populations will add equal numbers of individuals per unit time.
\correctchoice the large population will add more individuals per unit time.  
\choice the number of individuals added will depend on the value of N. 
\choice the small population will add more individuals per unit time. 
\end{choices}

\question
Symbiosis is a longer-term ecological interaction

\begin{choices}
\choice  where one species benefits and the other species neither benefits nor is harmed. 
\choice where one species benefits and another species is harmed 
\choice where two or more species partition resources. 
\correctchoice between two or more species. 
\choice where both species benefit. 
\end{choices}

\question
In ecosystems, why is the term cycling used to describe the transfer of chemical elements and nutrients, whereas the term flow is used for the transfer of energy?

\begin{choices}
\choice Chemical elements are cycled into ecosystems from other ecosystems, but energy constantly flows within the ecosystem.
\choice Energy flows from one ecosystem to other ecosystems, but chemical elements constantly cycle within an ecosystem.
\choice Both chemical elements and energy flow within an ecosystem but can cycle to other ecosystems.
\correctchoice Chemical elements are recycled within ecosystems, but energy flows through and out of ecosystems.
\choice Both chemical elements and energy are recycled but can flow to other ecosystems.
\end{choices}

\question
The basic model of metapopulations views the number of suitable habitat patches occupied at a given point in time as a balance between

\begin{choices}
\choice resource quality and degree of isolation. 
\choice dispersal and birth rate. 
\choice birth rate and local extinction. 
\correctchoice recolonization and local extinctions. 
\choice recolonization and degree of isolation. 
\end{choices}

\newpage

\question
A patch in a population has the following four values. 

\qquad\begin{tabular}{@{}ll@{}}
$b = 0.04$ & $i = 0.06$\\
$d = 0.08$ & $e = 0.04$\\
\end{tabular}

Based only on these values (ignoring habitat quality, patch size, etc.), the patch is most likely to be a 

\begin{choices}
\choice  source. 
\correctchoice  sink.
\end{choices}

\question
The attribute that most influence population density and rates of immigration and emigration among habitat patches in a metapopulation are

\begin{choices}
\choice habitat quality, size of sink patches and isolation of source patches. 
\choice patch isolation, habitat quality, and size of source patches. 
\choice patch size, habitat quality, and isolation of sink patches. 
\correctchoice patch size, patch isolation and habitat quality. 
\choice patch size, patch isolation and species diversity. 
\end{choices}

\question
Assume that you performed an \textsc{anova}. The final probability was $p = 0.035$. Based on this information, which of the following is true?

\begin{choices}
\choice You would accept the alternative hypothesis that the samples are from the same statistical population. 
\correctchoice You would reject the null hypothesis that the samples are from the same statistical population. 
\choice You would reject the alternative hypothesis that the samples are from different statistical populations.
 
\choice You would reject the alternative hypothesis that the samples are from the same statistical population. 
\choice You would accept the null hypothesis that the samples are from different statistical populations. 
\choice You would accept the null hypothesis that the samples are from the same statistical population. 
\end{choices}

\question
A field containing 2000 kg of plant material can support about \rule{1in}{0.4pt} of tertiary consumers.

\begin{choices}
\correctchoice 2 kg 
\choice 20 kg 
\choice 10 kg 
\choice 10\% 
\choice 90\% 
\choice 200 kg 
\end{choices}

\question
The science of naming a new species is called

\begin{choices}
\correctchoice  taxonomy. 
\choice systematics. 
\choice binomial nomenclature. 
\choice phylogeny. 
\choice classification. 
\end{choices}

\newpage

\question
A secondary consumer has high abundance. In a trophic cascade, which of the following is probably \emph{not} true.
\begin{choices}
\choice Actually, all of the above are true. 
\choice The different trophic levels will have alternating high and low abundances. 
\choice The tertiary consumer will have low abundance. 
\correctchoice The primary producer will have low abundance. 
\choice The primary consumer will have low abundance. 
\end{choices}

\question
Which block in the figure below shows the greatest species diversity? (Different shades, including white, represent different species. The small squares indicate relative abundance of individuals present.)

\begin{multicols}{3}
\begin{choices}
\choice \includegraphics{diversity_mcD} 
\correctchoice \includegraphics{diversity_mcE} 
\choice \includegraphics{diversity_mcA} 
\choice \includegraphics{diversity_mcC} 
\choice \includegraphics{diversity_mcB} 
\end{choices}
\end{multicols}

\question
Competition for limited resources may influence gene flow by causing

\begin{choices}
\choice  an increased death rate, thereby reducing gene flow.
\choice  some individuals to evolve a new trait to be better competitors.
\choice  some individuals to partition resources, reducing intraspecific competition.
\correctchoice some individuals to migrate to another population.
\choice  some individuals to move to a different source patch in the same population.
\end{choices}

\question
The goal of hierarchical classification is to

\begin{choices}
\correctchoice create monophyletic groups that represent evolutionary history. 
\choice provide a testable hypothesis about organisms with similar structures. 
\choice use a phylogenetic tree to accurately name taxa. 
\choice represent all life on earth with only eight categories, from domain to species. 
\choice provide an easy to memorize arrangement of life on earth. 
\end{choices}

\newpage

\question
While observing a population of lizards during one week, you notice that the number of adults has increased and is higher at the end of the week than at the start. The number of juveniles is unchanged. The most likely explanation for such an observation would be 

\begin{choices}
\choice decreased emigration.  
\choice decreased death rate. 
\choice decreased immigration. 
\choice increased birth rate.  
\correctchoice increased immigration.  
\choice increased emigration.  
\end{choices}

\question
A change in the body location where two \textit{Hox} genes are expressed caused swimming appendages to evolve into feeding appendages in crabs and other crustaceans. This type of genetic change is an example of

\begin{choices}
\choice an adaptation to living in an quiet aquatic environment where fewer swimming appendages are needed. 
\choice heterochrony that caused a difference in developmental timing. 
\correctchoice a change in a toolkit gene that altered the arrangement of parts in the body plan. 
\choice paedomorphosis that caused the adult to retain the juvenile feeding appendage. 
\choice allometry that caused the swimming appendages to grow more slowly so they could be used for feeding. 
\end{choices}

\question
Often the growth cycle of one population has an effect on the cycle of another. As moose populations increase, for example, wolf populations also increase. If you consider the logistic equation for the wolf population, which of the factors accounts for the effect of the moose population on the size of the wolf population?

\begin{choices}
\choice $rN$
\choice $\Delta N$
\choice $N$
\choice $\Delta t$
\correctchoice $K$
\choice $r$
\end{choices}

\question
Given the following values and assuming exponential population growth, what will be the total population size after one generation has elapsed? Round to nearest whole individual.\vspace{1ex}

\begin{minipage}{1em}
\end{minipage}%
\hspace{2em}%
\begin{minipage}{1.5in}
$b = 0.09$ \\
$d = 0.12$ \\
$N = 18,657$\\
$K = 15,000$\\
\end{minipage}

\begin{choices}
\correctchoice $18,097$
\choice $3,657$
\choice $19,217$
\choice $15,000$
\choice $560$
\end{choices}

\newpage

\question
Which of the following statements is consistent with the principle of competitive exclusion?  

\begin{choices}
\choice The density of one competing species will have a positive impact on the population growth of the other competing species. 
\choice Evolution tends to increase competition between related species.   
\choice Two species with the same fundamental niche will exclude other competing species.   
\choice Bird species generally do not compete for nesting sites.  
\correctchoice Even a slight reproductive advantage will eventually lead to the elimination of the less well adapted of two competing species. 
\end{choices}

\question
In a mature forest of oak, maple, and hickory trees, a disease causes a reduction in the number of acorns produced by oak trees. Which of the following would \emph{least} likely be a \emph{direct} result of this? 

\begin{choices}
\choice There might be fewer squirrels because they feed on acorns produced by oak trees. 
\correctchoice There might be a decrease in the number of maple seeds as the disease spreads to other trees in the forest. 
\choice  There might be fewer mice and seed-eating birds because squirrels that normally eat acorns would eat more seeds and compete with the mice and birds.
\choice There might be an increase in the number of hickory trees because the competition between hickory nuts and acorns for germination sites would be reduced or eliminated. 
\choice There might be fewer owls because they feed on baby squirrels, mice, and young seed-eating birds, whose populations would be reduced. 
\end{choices}

\question
Which of the following is most likely to contribute to density-dependent regulation of populations?

\begin{choices}
\correctchoice intraspecific competition for nutrients
\choice earthquakes
\choice the removal of toxic waste by decomposers
\choice fires
\choice floods
\end{choices}

\question
White-breasted nuthatches and brown creepers both eat insects that hide in the furrows of bark in hardwood trees. The brown creeper searches for insects by hunting from the bottom of the tree trunk toward the top, whereas the white-breasted nuthatch searches from the top of the trunk down. Their feeding areas rarely overlap on the tree trunk. These hunting behaviors best illustrate which of the following ecological concepts?

\begin{choices}
\choice competitive exclusion
\correctchoice resource partitioning
\choice keystone species
\choice interspecific competition
\choice intraspecific competition
\end{choices}

\newpage

\question
In the experiment shown in the figure below, \textit{Balanus} was removed from the habitat shown on the left. Which of the following statements is a valid conclusion of this experiment?\\
\includegraphics[width=0.9\textwidth]{balanus.png}

\begin{choices}
\correctchoice The removal of \textit{Balanus} shows that the realized niche of \textit{Chthamalus} is smaller than its fundamental niche.
\choice These two species of barnacle do not show competitive exclusion.
\choice \textit{Balanus} can survive only in the lower intertidal zone because it is unable to resist dehydration.
\choice Balanus is inferior to \textit{Chthamalus} in competing for space on rocks lower in the intertidal zone. 
\choice If \textit{Chthamalus} were removed, \textit{Balanus}'s fundamental niche would become larger.
\end{choices}


\question
Which of the following is the correct and complete arrangement of the hierarchical classification. 
\begin{choices}
\choice  Domain, Kingdom, Phylum, Family, Class, Species, Genus.
\choice  Domain, Phylum, Family, Class, Order, Genus, Species.
\choice  Domain, Kingdom, Phylum, Order, Class, Genus, Family, Species.
\correctchoice  Domain, Kingdom, Phylum, Class, Order, Family, Genus, Species.
\choice  Domain, Kingdom, Phylum, Order, Class, Family, Genus, Species.
\end{choices}

\question
Based on the hypothetical food web below, choose one letter that represents an organism that could be both a secondary and tertiary consumer? The arrows point from lower to higher trophic levels. More than one answer is possible. Circle only one answer.

\begin{multicols}{3} 
\begin{choices}
\correctchoice A
\correctchoice E
\choice B
\choice D
\choice C
\end{choices}
\columnbreak
\includegraphics{foodweb.png}
\columnbreak
\end{multicols}
\vspace{-1\baselineskip}

